\subsection{Strongly Gauduchon metrics}
\label{sec:lee-gauduchon-strong}

Vanishing of the Lee--Gauduchon class is weaker than the strongly Gauduchon condition. We give an explicit example, following the distinction in \cite[Remarks 5.6--5.8]{OV25LeeGauduchon}. It also shows why the image of $j_X$ must not be replaced by all de Rham-exact Bott--Chern classes.

\begin{definition}
A Hermitian metric $\gamma$ is \emph{strongly Gauduchon} if $\gamma^{n-1}$ is the $(n-1,n-1)$-component of a smooth closed real $(2n-2)$-form.
\end{definition}

Writing that form as
\[
 \Theta=\gamma^{n-1}+\beta+\overline\beta,
 \qquad \beta\text{ of type }(n,n-2),
\]
shows that the definition is equivalent to $\partial(\gamma^{n-1})$ being $\bar\partial$-exact: its value is $-\bar\partial\beta$. Conversely, this equation and its conjugate give $\mathrm{d}\Theta=0$, because the remaining types exceed dimension $n$. Applying $\bar\partial$ also shows that $\gamma$ is Gauduchon. Every balanced metric is strongly Gauduchon by taking $\Theta=\gamma^{n-1}$.

\begin{proposition}\label{prop:lee-strong-obstruction}
A strongly Gauduchon metric has zero Lee--Gauduchon class. If a connected compact complex manifold admits a nonzero positive de Rham-exact $(1,1)$-current, it has no strongly Gauduchon metric, and hence no balanced metric.
\end{proposition}
\begin{proof}
Put $\Gamma=\gamma^{n-1}$ and use $\beta$ as above. In our normalization,
\[
 \dc\Gamma
 =\frac{\mathrm i}{4\pi}
       (\bar\partial\beta-\partial\overline\beta)
 =\mathrm{d}\!\left(\frac{\mathrm i}{4\pi}
                              (\beta-\overline\beta)\right).
\]
Thus $[\dc\Gamma]=0$.

\Needspace{8\baselineskip}
For the obstruction, suppose $T=\mathrm{d}S\geq0$ is nonzero and $\Theta$ is a closed real form with strictly positive $(n-1,n-1)$-part $\Gamma$. Type considerations and distributional Stokes give
\[
 \int_XT\wedge\Gamma
 =\int_XT\wedge\Theta=0.
\]
The first integral is strictly positive because $T\ne0$ and $\Gamma>0$, a contradiction. The argument uses the full current and a smooth test form; it does not involve a non-pluripolar product.
\end{proof}

\begin{example}\label{ex:lee-calabi-eckmann}
There is a connected compact complex threefold $X$ with
\[
 \operatorname{LG}(X)=\{0\},\qquad V_X=K_X=W_X=0,
\]
but with no strongly Gauduchon metric. Moreover, $X$ has a nonzero pseudoeffective Bott--Chern class whose de Rham image is zero and which is not in the image of $j_X$.
\end{example}
\begin{proof}
Consider the holomorphic action of the additive group $\C$ on
\[
 \widetilde X=(\C^2\setminus\{0\})\times(\C^2\setminus\{0\}),
 \qquad
 t\cdot(z,w)=(e^t z,e^{\mathrm i t}w).
\]
The action is free. Indeed, $e^t=e^{\mathrm i t}=1$ forces $t=0$. It is proper: if both $(z,w)$ and $t\cdot(z,w)$ range in compact subsets of $\widetilde X$, the norms of all four vector components are bounded above and away from zero. Writing $t=a+\mathrm i b$, the ratios of the norms bound $e^a$ and $e^{-b}$ above and away from zero. Thus $a$ and $b$ stay bounded, which proves properness. The free proper holomorphic action has local holomorphic slices and a Hausdorff complex quotient $X$ of dimension three.

Every orbit meets $S^3\times S^3$ exactly once. The required group element is determined by
\[
 a=-\log|z|,\qquad b=\log|w|.
\]
These expressions are smooth, so the slice identifies the quotient smoothly with $S^3\times S^3$. In particular $X$ is compact and connected, and
\[
 H^1_{\mathrm{dR}}(X,\R)=H^2_{\mathrm{dR}}(X,\R)
 =H^5_{\mathrm{dR}}(X,\R)=0.
\]
It follows directly that $V_X=K_X=W_X=0$ and $\operatorname{LG}(X)=\{0\}$: the latter cone is nonempty by Gauduchon's theorem and lies in the zero group $H^5_{\mathrm{dR}}(X,\R)$.

\Needspace{16\baselineskip}
Projectivizing both vector components descends to a holomorphic submersion
\[
 p:X\longrightarrow\mathbf P^1\times\mathbf P^1.
\]
Its differential is surjective because the product projectivization on $\widetilde X$ is a submersion and the group orbits lie in its fibers. Let $\eta$ be the sum of the Fubini--Study forms on the two factors. The smooth form
\[
 \alpha=p^*\eta
\]
is closed, semipositive and of rank two everywhere. It is nonzero and de Rham exact, since $H^2_{\mathrm{dR}}(X,\R)=0$. By \cref{prop:lee-strong-obstruction}, it obstructs a strongly Gauduchon metric.

Its Bott--Chern class is nonzero. Otherwise $\alpha=\ddc f$ for a smooth real function $f$, and integration against any Gauduchon form would give zero, contradicting positivity and $\alpha\ne0$. The class is pseudoeffective because it has the positive representative $\alpha$. Finally $j_X$ has zero domain, so this class is not in its image.
\end{proof}

The example distinguishes two linear obstructions. The Lee--Gauduchon cone detects positive classes obtained from closed one-forms by $j_X$. Strongly Gauduchon metrics are obstructed by every nonzero positive de Rham-exact $(1,1)$-current. The latter family can be larger, even when the first cohomology vanishes. Thus the bimeromorphic invariant above is useful without being a complete test for balanced or strongly Gauduchon geometry.
